% ============================================================
% §03b — GIS: GERT Interpolation Syntax (Normative)
% Author:  Edith — Spec Editor
% Created: 2026-06-05
% Source of truth: design/gert/grammar/gis.ebnf (v1.0.0-draft)
% ============================================================
%
% RFC 2119 normative language used throughout:
%   MUST / MUST NOT  — absolute requirements for conforming implementations.
%   SHOULD / SHOULD NOT — recommended but not mandatory.
%   MAY              — optional; permitted but not required.
%
% Grammar production citations take the form: gis.ebnf §<section>.<subsection>.
% Decision ledger citations take the form: decisions.md "<heading>" <key>.

% Local definition — remove if \TODO is promoted to the main preamble.
\providecommand{\TODO}[1]{\textbf{\textcolor{red}{[TODO:~#1]}}}

% ─────────────────────────────────────────────────────────────────────────────
\chapter{GIS\,---\,GERT Interpolation Syntax}
\label{sec:gis}
% ─────────────────────────────────────────────────────────────────────────────

\begin{quote}
\noindent\textbf{Informative introduction.}
GIS (GERT Interpolation Syntax) is the GERT-native template language for
embedding run-time values into string-typed runbook fields.
Wherever a field carries dynamic content — step \texttt{title},
command \texttt{args}, \texttt{stdin}, tool \texttt{argv} elements,
\texttt{display.content}, assertion \texttt{subject} and \texttt{expected}
operands, artifact paths, \texttt{event.id}, and \texttt{include.with}
values — the field value is treated as a GIS template and evaluated
immediately before the containing step executes.

A GIS template is a Unicode string that may contain zero or more
interpolation blocks of the form \texttt{\$\{}\textit{expr}\texttt{\}},
where \textit{expr} is a GXL expression (\S\ref{sec:gxl}).
The most common use is a bare GDP path
(e.g., \texttt{\$\{hostname\}}, \texttt{\$\{config.region\}}),
but the full GXL grammar is available inside every interpolation block.

GIS is the only supported interpolation surface.
This chapter is the authoritative normative statement of the
GERT Portable JSON Value Model (PJVM) used by GIS, GXL, and GCP.
\end{quote}

% ─────────────────────────────────────────────────────────────────────────────
\section{Normative Reference to the Grammar}
\label{sec:gis:normref}
% ─────────────────────────────────────────────────────────────────────────────

The file \texttt{design/gert/grammar/gis.ebnf} (version~1.0.0-draft) is the
authoritative grammar for GIS\@.  This chapter is normative prose \emph{about}
that grammar; it does not supersede it.  Where this prose and the grammar
conflict, \textbf{the grammar wins} and the prose is in error.  Conforming
implementations MUST derive their template parsers from \texttt{gis.ebnf},
not from this chapter alone.

GIS is a \emph{two-level} grammar: the outer level is defined in
\texttt{gis.ebnf} and processes the template string; the inner level is the
GXL grammar (\texttt{gxl.ebnf}) that processes each interpolation expression.
Implementations MUST apply the two parsers in sequence and MUST NOT conflate
their token streams.

(\texttt{gis.ebnf} \S1, \S2.)

% ─────────────────────────────────────────────────────────────────────────────
\section{Lexical Structure}
\label{sec:gis:lexical}
% ─────────────────────────────────────────────────────────────────────────────

The GIS lexer operates over the raw string value of the YAML field,
after YAML parsing has resolved all YAML-level escaping and produced a
plain Unicode string.
The GIS lexer then partitions that string into a sequence of segments.
All productions in this section correspond to \texttt{gis.ebnf} \S2.

% ─────────────────────────────────────────────────────────────────────────────
\subsection{Template String Structure}
\label{subsec:gis:template-string}

The top-level GIS production is:

\begin{minted}{text}
TemplateString = { RawSegment | EscapeSequence | Interpolation } ;
\end{minted}

A \texttt{TemplateString} containing zero \texttt{Interpolation} blocks MUST
be returned as-is, without alteration.
No error is raised for a template string that contains only raw characters.

(\texttt{gis.ebnf} \S2.1.)

% ─────────────────────────────────────────────────────────────────────────────
\subsection{Raw Segments}
\label{subsec:gis:raw}

\begin{minted}{text}
RawSegment = RAW_CHAR { RAW_CHAR } ;
RAW_CHAR   = (* any Unicode code point except '\' and '$' *)
           | '$' (* ONLY when NOT immediately followed by '{' *)
           ;
\end{minted}

A bare \texttt{\$} character that is \emph{not} immediately followed by
\texttt{\{} is treated as a literal dollar sign and passed through verbatim.
This permits shell-style environment-variable references
(e.g., \texttt{\$PATH}) to coexist with GIS interpolations in
\texttt{argv} values, provided the \texttt{\$} is not followed by \texttt{\{}.

\begin{quote}
\noindent\textit{Non-normative note.}
Open issue OI-GIS-02 (\texttt{gis.ebnf} \S6) identifies a class of tool
definitions where both shell-style \texttt{\$\{ENV\_VAR\}} and GIS
\texttt{\$\{variable\}} appear in the same \texttt{argv} string.
Under the current grammar ALL \texttt{\$\{...\}} sequences are GIS
interpolations; a literal \texttt{\$\{ENV\_VAR\}} intended for the OS
process environment MUST be written as \texttt{\textbackslash\$\{ENV\_VAR\}}.
\end{quote}

(\texttt{gis.ebnf} \S2.2.)

% ─────────────────────────────────────────────────────────────────────────────
\subsection{Escape Sequences}
\label{subsec:gis:escapes}

\begin{minted}{text}
EscapeSequence = '\' ESCAPE_TARGET ;
ESCAPE_TARGET  = '$' '{'   (* '\${' → literal '${' in output    *)
               | '\'        (* '\\' → literal '\'                *)
               | 'n'        (* '\n' → U+000A LINE FEED           *)
               | 'r'        (* '\r' → U+000D CARRIAGE RETURN     *)
               | 't'        (* '\t' → U+0009 CHARACTER TABULATION*)
               ;
\end{minted}

The sequence \texttt{\textbackslash\$\{} is the \textbf{canonical escape}
for a literal dollar-brace pair \texttt{\$\{} in template output.
This is the single backslash followed by \texttt{\$\{} (OI-GIS-01, ratified;
decisions.md ``GXL Stream~A open issue ratifications'').

\noindent The recognized escape sequences and their output characters are:

\begin{center}
\begin{tabular}{lll}
\hline
\textbf{Sequence} & \textbf{Output} & \textbf{Description} \\
\hline
\texttt{\textbackslash\$\{} & \texttt{\$\{}        & Literal dollar-brace (canonical escape; OI-GIS-01) \\
\texttt{\textbackslash\textbackslash} & \texttt{\textbackslash} & Literal backslash \\
\texttt{\textbackslash n}            & U+000A              & Line feed \\
\texttt{\textbackslash r}            & U+000D              & Carriage return \\
\texttt{\textbackslash t}            & U+0009              & Horizontal tab \\
\hline
\end{tabular}
\end{center}

\noindent\textbf{Lenient unknown-escape policy.}
A backslash followed by any character \texttt{X} not in the list above
(i.e., not \texttt{\$\{}, \texttt{\textbackslash}, \texttt{n}, \texttt{r},
or \texttt{t}) is passed through as the two-character literal
\texttt{\textbackslash X} and MUST generate a \texttt{GIS-PARSE-002}
diagnostic (a \emph{warning}, not a hard error; see \S\ref{sec:gis:errors}).
This lenient behavior differs intentionally from GXL
(\texttt{GXL-PARSE-004}, which is a hard error) because GIS template
strings frequently arrive from YAML values that have already undergone
YAML-level escape processing; strict rejection would break runbooks that
contain literal backslashes in tool argument strings.

(\texttt{gis.ebnf} \S2.3.)

% ─────────────────────────────────────────────────────────────────────────────
\subsection{Interpolation Blocks}
\label{subsec:gis:interpolation}

\begin{minted}{text}
Interpolation = '$' '{' GISExpression '}' ;
\end{minted}

The two-character sequence \texttt{\$\{} is the start delimiter of an
interpolation block.
The block ends at the first \texttt{\}} that is not inside a nested
grouping within the \texttt{GISExpression}
(see \S\ref{sec:gis:expression} for nesting rules).

An unclosed \texttt{\$\{} — i.e., a \texttt{\$\{} that has no matching
\texttt{\}} before the end of the template string — MUST cause the
entire template to be rejected with error \texttt{GIS-PARSE-001}.

An unmatched \texttt{\}} outside an interpolation block is a literal
\texttt{\}} and is passed through verbatim.
An unmatched \texttt{\{} outside an interpolation block is likewise
a literal \texttt{\{}.
Only the two-character sequence \texttt{\$\{} begins an interpolation.

(\texttt{gis.ebnf} \S2.4.)

% ─────────────────────────────────────────────────────────────────────────────
\section{Embedded Expression}
\label{sec:gis:expression}
% ─────────────────────────────────────────────────────────────────────────────

The content of each interpolation block is a GXL expression:

\begin{minted}{text}
GISExpression = Expression ;
\end{minted}

\texttt{Expression} is the top-level non-terminal of the GXL syntactic
grammar (\texttt{gxl.ebnf} \S3; \S\ref{sec:gxl}).
Conforming implementations MUST support the full GXL grammar inside
\texttt{\$\{...\}}, not a restricted subset.
The GXL lexer handles string tokenization within the expression, which
means the outer GIS parser MUST delegate to the GXL parser's
``end of expression'' signal to locate the closing \texttt{\}};
naive character-scanning of \texttt{\}} will mis-parse expressions
containing string literals with embedded \texttt{\}}.
(See \texttt{gis.ebnf} \S6, OI-GIS-04.)

\noindent\textbf{In practice}, the overwhelming majority of interpolation
sites use a bare GDP path (\S\ref{subsec:gxl:gdp}), which is a valid
GXL sub-expression:

\begin{minted}{text}
${hostname}              -- top-level variable (bare IDENT)
${config.region}         -- field access on captured object
${items[0].name}         -- index into captured array, then field
\end{minted}

Full GXL expressions are equally valid:

\begin{minted}{text}
${str.toLower(env)}      -- stdlib call; result coerced to string
${count + 1}             -- arithmetic; result coerced to string
${status == "ok"}        -- comparison; "true" or "false" in output
\end{minted}

The result of the GXL expression is coerced to a string value before
insertion into the template output, according to the rules in
\S\ref{sec:gis:coercion}.

\noindent\textbf{Empty interpolation.}
The block \texttt{\$\{\}} (dollar-brace with no expression) MUST be
rejected at parse time with error \texttt{GIS-PARSE-004}.
There is no implicit expression for an empty block.

(\texttt{gis.ebnf} \S2.4, \S3.)

% ─────────────────────────────────────────────────────────────────────────────
\section{GERT Portable JSON Value Model}
\label{sec:gis:portable-json}
% ─────────────────────────────────────────────────────────────────────────────

\begin{quote}
\noindent\textbf{Authoritative definition.}
This section is the single normative home of the GERT Portable JSON Value
Model (PJVM).
All other GERT specification chapters that reference PJVM types
(\S\ref{sec:gxl}, \S\ref{sec:gcp}) MUST treat this section as the
authoritative definition and MUST NOT duplicate or amend it.
\end{quote}

The GERT Portable JSON Value Model (PJVM) is the runtime type system
shared by GXL (\S\ref{sec:gxl}), GIS, and GCP (\S\ref{sec:gcp}).
It is derived from RFC~8259 (JSON) with the explicit constraints
enumerated in this section.
The design goal of PJVM is cross-runtime fidelity: a value serialized
in the Go reference runtime and deserialized in the C\# runtime (or any
future runtime) MUST have an identical PJVM representation with no
runtime-specific surprises.

(\texttt{gis.ebnf} \S4; OQ3, resolved.)

% ─────────────────────────────────────────────────────────────────────────────
\subsection{PJVM Types}
\label{subsec:gis:pjvm-types}

PJVM defines exactly six value types:

\begin{center}
\begin{tabular}{lp{9.5cm}}
\hline
\textbf{Type} & \textbf{Definition} \\
\hline
\texttt{null}    & The absence of a value, corresponding to JSON \texttt{null}.
                   \texttt{null} is a distinct type; it is not the same as
                   an empty string or the number zero. \\
\texttt{boolean} & One of the two Boolean values: \texttt{true} or
                   \texttt{false}, corresponding to JSON Boolean literals.
                   GXL type-checking error messages use the abbreviation
                   \texttt{bool} for this type. \\
\texttt{number}  & A finite IEEE~754 double-precision float64 value.
                   \texttt{NaN} and \texttt{Infinity} are \textbf{not}
                   members of this type; see the exclusion table below. \\
\texttt{string}  & An immutable, UTF-8-encoded sequence of Unicode code
                   points, corresponding to a JSON string value.
                   No length limit is mandated by this specification;
                   implementations SHOULD document any practical limit. \\
\texttt{array}   & An ordered, homogeneous-or-heterogeneous sequence of
                   PJVM values, corresponding to a JSON array.
                   Elements are zero-indexed. \\
\texttt{object}  & An unordered collection of string-keyed PJVM values,
                   corresponding to a JSON object.
                   Key collisions in a source document are resolved by
                   retaining the last occurrence (RFC~8259 \S4 behavior). \\
\hline
\end{tabular}
\end{center}

(\texttt{gis.ebnf} \S4.1.)

% ─────────────────────────────────────────────────────────────────────────────
\subsection{PJVM Exclusions}
\label{subsec:gis:pjvm-exclusions}

The following value types are explicitly excluded from PJVM and MUST NOT
appear in the run variable context map:

\begin{center}
\begin{tabular}{lp{9.5cm}}
\hline
\textbf{Excluded Type} & \textbf{Rationale and Mapping} \\
\hline
\texttt{NaN}       & IEEE~754 not-a-number.
                     A capture that would yield \texttt{NaN} MUST
                     substitute \texttt{null} and raise
                     \texttt{GIS-TYPE-001} if the value reaches
                     a GIS interpolation site. \\
\texttt{Infinity}  & IEEE~754 positive or negative infinity.
                     Same policy as \texttt{NaN}. \\
\texttt{undefined} & JavaScript \texttt{undefined}.
                     MUST be mapped to \texttt{null} at ingestion. \\
\texttt{date}      & No date type.  Dates are represented as strings
                     in a documented format (ISO~8601 RECOMMENDED). \\
\texttt{binary}    & No blob or buffer type.
                     Binary data MUST be base64-encoded to a string. \\
\texttt{int64}     & No distinct integer type.
                     All numbers are \texttt{float64}; integers are
                     representable exactly up to $2^{53}$. \\
\texttt{decimal}   & No arbitrary-precision decimal.
                     All numbers are \texttt{float64}. \\
\hline
\end{tabular}
\end{center}

(\texttt{gis.ebnf} \S4.1.)

% ─────────────────────────────────────────────────────────────────────────────
\subsection{Unresolved Variables}
\label{subsec:gis:unresolved}

If the GXL expression inside \texttt{\$\{...\}} references a variable
that is not present in the run variable context map (error
\texttt{GXL-PATH-001}; \S\ref{sec:gxl:semantics:paths}), the GIS evaluator
MUST propagate the failure as \texttt{GIS-PATH-MISSING} and MUST abort
template rendering.

Silent empty-string substitution for missing variables is \textbf{forbidden}.
Every unresolved reference is a hard error.
Authors MUST use \texttt{when:} guards (exploiting GXL short-circuit
evaluation; \S\ref{sec:gxl:semantics:short-circuit}) or
\texttt{capture.default:} values (\S\ref{sec:gcp:default}) to handle
optional references.

(\texttt{gis.ebnf} \S4.3.)

% ─────────────────────────────────────────────────────────────────────────────
\section{Type Coercion to String}
\label{sec:gis:coercion}
% ─────────────────────────────────────────────────────────────────────────────

When a GXL expression inside \texttt{\$\{...\}} resolves to a PJVM value,
that value MUST be converted to a string before insertion into the template
output.
The coercion rules are normative and are exhaustive:
no PJVM type is an error to interpolate, except as explicitly noted below.

\begin{center}
\begin{tabular}{lp{9.5cm}}
\hline
\textbf{PJVM Type} & \textbf{Coercion to String} \\
\hline
\texttt{null}    & The four-character string \texttt{null}.
                   The empty string is not a valid coercion of \texttt{null}. \\
\texttt{boolean} & \texttt{true} or \texttt{false} (the keyword spelling),
                   matching the GXL literal syntax. \\
\texttt{number}  & Shortest round-trip decimal representation satisfying
                   \mbox{$\mathrm{parse}(\mathrm{repr}(x)) = x$}.
                   Full rules in \S\ref{subsec:gis:coercion-number} below. \\
\texttt{string}  & The string value, verbatim.
                   No quoting, no escaping is added. \\
\texttt{array}   & Canonical JSON serialization: \texttt{[} followed by
                   comma-separated element representations followed by
                   \texttt{]}.  No extra whitespace.
                   Array elements are recursively coerced to their canonical
                   JSON forms (not to strings first).
                   Example: \texttt{[1,"a",true]} \\
\texttt{object}  & Canonical JSON serialization with
                   \textbf{lexicographically sorted keys}
                   (Unicode code point order, ascending).
                   Format: \texttt{\{} comma-separated \texttt{"key":value}
                   pairs \texttt{\}}.  No extra whitespace.
                   Example: \texttt{\{"a":2,"z":1\}} \\
\hline
\end{tabular}
\end{center}

\noindent\textbf{Rationale for sorted object keys.}
Canonical serialization with sorted keys ensures that interpolating a
captured object variable (e.g., \texttt{\$\{config\}}) produces the same
string regardless of which runtime generated the object,
making audit trail comparisons deterministic and reliable.

(\texttt{gis.ebnf} \S4.2.)

% ─────────────────────────────────────────────────────────────────────────────
\subsection{Number Coercion Rules}
\label{subsec:gis:coercion-number}

Number-to-string coercion MUST produce the shortest decimal representation
that round-trips through standard \texttt{float64} parsing
(\texttt{gis.ebnf} \S4.2).
Implementations SHOULD use ECMAScript's \textsc{Number::toString} algorithm
(ECMA-262 \S7.1.12.1) or Go's \texttt{strconv.FormatFloat(x, 'g', -1, 64)},
both of which satisfy the round-trip requirement.
The following rules are normative:

\begin{itemize}
  \item \textbf{Integer-valued floats}: no decimal point is emitted.
    \texttt{42.0} coerces to \texttt{"42"}, not \texttt{"42.0"}.

  \item \textbf{Non-integer floats}: the minimum number of decimal digits
    required for round-trip fidelity is used.
    \texttt{3.14} coerces to \texttt{"3.14"}.

  \item \textbf{Large integers}: values below $10^{21}$ are represented
    without exponential notation.
    \texttt{1e20} coerces to \texttt{"100000000000000000000"}.
    Values at or above $10^{21}$ use exponential notation with a
    \texttt{+} in the exponent: \texttt{1e21} coerces to
    \texttt{"1e+21"}.

  \item \textbf{Negative zero}: \texttt{-0.0} coerces to \texttt{"0"},
    not \texttt{"-0"}.

  \item \textbf{NaN and Infinity}: these values are not members of PJVM
    (see \S\ref{subsec:gis:pjvm-exclusions}).
    A number value of \texttt{NaN} or \texttt{Infinity} encountered at
    interpolation time MUST cause the template rendering to fail with
    \texttt{GIS-TYPE-001}.
    The runtime SHOULD have already substituted \texttt{null} at capture
    time; reaching interpolation with a non-finite value indicates a
    runtime defect.
\end{itemize}

(\texttt{gis.ebnf} \S4.2.)

% ─────────────────────────────────────────────────────────────────────────────
\section{Error Class Catalog}
\label{sec:gis:errors}
% ─────────────────────────────────────────────────────────────────────────────

GIS defines two timing categories for errors, consistent with GXL
(\S\ref{sec:gxl:errors}):
\emph{parse errors} are raised at runbook load/plan time, before any step
executes;
\emph{evaluation errors} are raised at runtime when a template is rendered
for a specific step.
A runbook containing a parse error MUST NOT enter execution.
An evaluation error during template rendering causes the affected step
to transition to the \texttt{failed} state.

\texttt{GIS-PARSE-002} is a \emph{warning}: it does not prevent execution.
All other GIS errors are \emph{hard errors}.

The master PLAN-* code catalog is owned by \S\ref{sec:parse-gate:error-codes}.
GIS parse errors are surfaced at the parse gate via \texttt{PLAN-001}
(unparseable expression wrapper); the specific GIS error code is included
in the \texttt{sub\_error} field of the \texttt{PLAN-001} diagnostic.

Table~\ref{tab:gis-errors} enumerates every GIS error code.

\begin{table}[ht]
\centering
\begin{tabular}{p{3.5cm}p{2.0cm}p{5.0cm}p{2.5cm}}
\hline
\textbf{Code} & \textbf{Raised By} & \textbf{Description} & \textbf{Remediation} \\
\hline
\multicolumn{4}{l}{\textit{Parse errors (raised at load/plan time)}} \\
\hline
\texttt{GIS-PARSE-001}
  & GIS lexer
  & Unclosed interpolation block: \texttt{\$\{} appears in the template with no matching \texttt{\}}.  The entire template is rejected.
  & Close the \texttt{\$\{} with a matching \texttt{\}}. \\
\texttt{GIS-PARSE-002}
  & GIS lexer
  & Unknown escape sequence: \texttt{\textbackslash X} where \texttt{X} is not a recognized escape target (\S\ref{subsec:gis:escapes}).  \emph{Warning only}; \texttt{\textbackslash X} is passed through literally.
  & Use a recognized escape, or double the backslash: \texttt{\textbackslash\textbackslash X}. \\
\texttt{GIS-PARSE-003}
  & GIS parser
  & Invalid GXL expression inside \texttt{\$\{...\}}: the GXL parser raised a \texttt{GXL-PARSE-*} error.  The specific GXL code is included in the \texttt{GIS-PARSE-003} diagnostic.
  & Fix the flagged expression (see GXL error catalog, \S\ref{sec:gxl:errors}). \\
\texttt{GIS-PARSE-004}
  & GIS parser
  & Empty interpolation block: \texttt{\$\{\}} with no expression.
  & Provide a GXL expression inside the braces. \\
\hline
\multicolumn{4}{l}{\textit{Path errors (raised at evaluation time)}} \\
\hline
\texttt{GIS-PATH-MISSING}
  & GIS evaluator
  & The GXL expression inside \texttt{\$\{...\}} raised \texttt{GXL-PATH-001}: the referenced variable is not in the run variable context.  The template location and variable name are included in the diagnostic.
  & Ensure the variable is captured before this step runs, or guard the step with \texttt{when:}. \\
\hline
\multicolumn{4}{l}{\textit{Type errors (raised at evaluation time)}} \\
\hline
\texttt{GIS-TYPE-001}
  & GIS evaluator
  & The resolved value cannot be coerced to string: specifically, a \texttt{number} value that is \texttt{NaN} or \texttt{Infinity} (\S\ref{subsec:gis:coercion-number}).
  & Ensure the capture source cannot produce non-finite numbers; use a null-guard default. \\
\texttt{GIS-TYPE-002}
  & GIS evaluator
  & A GXL type error was propagated from the embedded expression (e.g., \texttt{GXL-TYPE-001}, \texttt{GXL-TYPE-002}, \texttt{GXL-TYPE-003}).  The specific GXL error code is included in the diagnostic.
  & Correct the expression type mismatch (see GXL error catalog, \S\ref{sec:gxl:errors}). \\
\hline
\end{tabular}
\caption{GIS error class catalog.
  Source: \texttt{gis.ebnf} \S5.}
\label{tab:gis-errors}
\end{table}

% ─────────────────────────────────────────────────────────────────────────────
\section{Optional Path Chaining}
\label{sec:gis:optional-chaining}
% ─────────────────────────────────────────────────────────────────────────────

GIS optional path chaining is the explicit opt-in syntax for rendering an
empty string when a path segment that the author has marked optional is
missing.  The governing decision is the ratified GIS optional-chaining entry
(
\href{run:../../../.squad/decisions/inbox/copilot-gis-optional-chaining-ratified.md}{\texttt{copilot-gis-optional-chaining-ratified.md}},
``Ratified Decisions''); it is consistent with the Phase~1 decision ledger's
hard-error default for plain GIS paths (decisions.md ``GXL Phase~1 Day~2 ---
GIS/GCP Specs'').  The semantic source prose is Barbara's proposal
\texttt{design/gert/proposals/gis-optional-chaining.md} \S2.

\begin{quote}
\noindent\textbf{Normative scope.}
This section specifies optional chaining only for GDP path resolution inside
GIS interpolation blocks.  It does not amend GXL boolean expressions, GCP
capture paths, or the PJVM coercion rules except where explicitly stated here.
\end{quote}

% ─────────────────────────────────────────────────────────────────────────────
\subsection{Syntax}
\label{subsec:gis:optional-chaining:syntax}

The optional path segment forms are \texttt{?.}\textit{IDENT} and
\texttt{?.[}\textit{N}\texttt{]}.  Conforming GIS implementations MUST accept
both forms inside a GDP path in a GIS interpolation expression, according to
the ratified EBNF delta:

\begin{minted}{text}
GDP = IDENT { PathSegment } ;
PathSegment = '.' IDENT
            | '?.' IDENT
            | '[' INTEGER ']'
            | '?.[' INTEGER ']'
            ;
\end{minted}

\noindent The full grammar lives in
\href{run:../grammar/gis.ebnf}{\texttt{design/gert/grammar/gis.ebnf}} and the
embedded GDP production is the same path form referenced from
\S\ref{subsec:gxl:gdp}; the subset shown here is only the ratified extension
surface.  The token \texttt{?.} MUST be recognized as an atomic optional-dot
segment separator, and \texttt{?.[} MUST be recognized as the start of an
optional bracket segment.  (Source: proposal \S1.1--\S1.2 and Appendix~A;
ratified decision rows ``\texttt{?.[N]} optional bracket indexing'' and
``Scope''.)

A GDP path MUST still begin with \texttt{IDENT}.  Therefore
\texttt{\$\{?.root\}} is syntactically illegal and MUST be rejected at parse
time: there is no grammar production for an optional root identifier.  The
legal spelling is an ordinary root identifier followed by an optional segment,
for example \texttt{\$\{root?.field\}}.  (Source: proposal \S1.1 and ratified
decision row ``\texttt{\$\{?.root\}} (optional on root identifier)''.)

% ─────────────────────────────────────────────────────────────────────────────
\subsection{Evaluation Semantics}
\label{subsec:gis:optional-chaining:semantics}

When a GIS interpolation expression evaluates a GDP path left-to-right, a miss
encountered by a \texttt{?.}\textit{IDENT} or \texttt{?.[}\textit{N}\texttt{]}
segment MUST stop evaluation of the entire remaining path tail and MUST resolve
the GDP expression to the empty string \texttt{""}.  No later segment in that
path is evaluated after the optional miss.  This is the ratified
JS/TS-compatible full-tail short-circuit rule, with GIS substituting
\texttt{""} for the JavaScript/TypeScript undefined result.  (Source: proposal
\S2.1; ratified decision rows ``Default value for missing optional path'' and
``Short-circuit semantics''.)

For optional chaining, the following conditions are missing and MUST trigger
that empty-string result when they are reached through an optional segment:

\begin{itemize}
  \item the root identifier is absent from the context map and the first path
    hop from that identifier is optional, as in \texttt{\$\{user?.name\}};
  \item the current value is \texttt{null};
  \item the current object lacks the requested key;
  \item the array index in \texttt{?.[}\textit{N}\texttt{]} is out of bounds.
\end{itemize}

\noindent Each condition above is normative.  (Source: proposal \S2.2;
ratified decision rows ``Default value for missing optional path'',
``Short-circuit semantics'', and ``\texttt{?.[N]} optional bracket indexing''.)

The following PJVM values are present values, not misses, and MUST NOT trigger
optional-chain defaulting: the empty string \texttt{""}, \texttt{false},
\texttt{0}, the empty array \texttt{[]}, and the empty object \texttt{\{\}}.
They continue through normal path evaluation and, when they become the final
interpolation result, are coerced according to \S\ref{sec:gis:coercion}.
(Source: proposal \S2.3; ratified decision row ``Default value for missing
optional path''.)

Plain GIS paths are unchanged.  A path written without optional segments, such
as \texttt{\$\{a.b.c\}}, MUST still fail with \texttt{GIS-PATH-MISSING} on any
missing identifier or key as specified in \S\ref{subsec:gis:unresolved}.  There
is no implicit empty-string substitution for ordinary \texttt{\$\{...\}}
interpolation.  (Source: \texttt{gis.ebnf} \S4.3; ratified decision row
``Hard-error default unchanged''.)

Mixed paths are legal, but non-optional segments remain mandatory.  If a
mandatory \texttt{.}\textit{IDENT} or \texttt{[}\textit{N}\texttt{]} segment
misses before any guarding optional segment applies, the GIS evaluator MUST
raise the ordinary hard error; it MUST NOT reinterpret that miss as optional.
(Source: proposal \S1.3 and \S2.4; ratified decision row ``Hard-error default
unchanged''.)

\begin{minted}{text}
Expression: ${a.b?.c.d}
Context:    { a: { b: { c: { d: "found" } } } }
Result:     "found"

Expression: ${a.b?.c.d}
Context:    { a: { b: {} } }
Result:     ""        -- ?.c misses; .d is not evaluated

Expression: ${a.b?.c.d}
Context:    { a: {} }
Result:     GIS-PATH-MISSING
            -- .b is mandatory; ?.c is never reached
\end{minted}

\noindent Optional bracket access follows the same full-tail rule:

\begin{minted}{text}
Expression: ${results?.[0]?.title}
Context:    { results: [] }
Result:     ""        -- ?.[0] misses; ?.title is not evaluated
\end{minted}

% ─────────────────────────────────────────────────────────────────────────────
\subsection{Composition with GXL Expressions}
\label{subsec:gis:optional-chaining:composition}

Optional chaining is resolved at the GDP path-resolution step.  It MUST NOT
propagate through enclosing GXL function calls.  For example, if
\texttt{user} is absent, \texttt{\$\{str.toLower(user?.name)\}} MUST first
resolve \texttt{user?.name} to \texttt{""} and then call
\texttt{str.toLower("")}.  The function receives an ordinary string argument;
it does not receive or observe an optional-chain marker.  (Source: proposal
\S3.2 and \S8 OQ-OC-4; ratified decision row ``Stdlib call propagation''.)

Optional chaining also composes with ordinary GXL expression evaluation after
the GDP has resolved.  If \texttt{user?.name} resolves to \texttt{""}, an
expression such as \texttt{\$\{user?.name + " suffix"\}} proceeds as string
concatenation with an empty left operand.  (Source: proposal \S3.2; ratified
decision row ``Stdlib call propagation''.)

% ─────────────────────────────────────────────────────────────────────────────
\subsection{Relationship to \texttt{capture.default:}}
\label{subsec:gis:optional-chaining:capture-default}

Optional chaining and \texttt{capture.default:} address different layers.
Optional chaining is a GIS rendering-time path rule for \texttt{\$\{...\}}
interpolation; \texttt{capture.default:} is a GCP capture-time rule for scalar
capture values (\S\ref{sec:gcp:default}).  Authors SHOULD use
\texttt{capture.default:} when the run variable itself should be populated with
a default at capture time, and SHOULD use \texttt{?.} only when the omission is
local to a rendered string.  (Source: proposal \S2.5 and \S3.1; ratified
decision rows ``Scope'' and ``\texttt{??} nullish-coalescing''.)

% ─────────────────────────────────────────────────────────────────────────────
\subsection{Examples}
\label{subsec:gis:optional-chaining:examples}

\noindent\textbf{Idiomatic: optional API response field.}

\begin{minted}{yaml}
steps:
  - tool: http.get
    args: "${api_base}/users/${user_id}"
    capture:
      http.body.profile: user_profile
  - display:
      content: |
        Name: ${user_profile.name}
        Email: ${user_profile.email}
        Company: ${user_profile?.company}
\end{minted}

\noindent If \texttt{company} is absent, the third line renders as
\texttt{Company: } rather than failing the template.  (Source: proposal \S5.1;
ratified decision row ``Default value for missing optional path''.)

\noindent\textbf{Idiomatic: optional captured subtree.}

\begin{minted}{yaml}
steps:
  - tool: api.call
    capture:
      http.body: response
  - display:
      content: "Request ID: ${response?.metadata?.request_id}"
\end{minted}

\noindent The \texttt{metadata} object and its \texttt{request\_id} member are
both optional for this rendered string.  (Source: proposal \S5.3; ratified
decision row ``Short-circuit semantics''.)

\noindent\textbf{Idiomatic: optional configuration section.}

\begin{minted}{yaml}
steps:
  - tool: deploy
    args: "--name ${service_name} --monitoring-endpoint ${config?.monitoring?.endpoint}"
\end{minted}

\noindent If \texttt{config.monitoring} is absent, the optional chain renders
the endpoint value as \texttt{""}; the receiving tool remains responsible for
whether an empty argument is acceptable.  (Source: proposal \S5.2; ratified
decision row ``Scope''.)

\noindent\textbf{Anti-pattern: hiding required-field bugs.}

\begin{minted}{yaml}
steps:
  - display:
      content: "Order total: ${order?.total}"
\end{minted}

\noindent If \texttt{order} is required at this point in the runbook, the
example above is an anti-pattern: it hides a missing capture, a misspelled
variable, or a broken step dependency.  The parser cannot prove the author's
intent here; use \texttt{?.} only after deciding that the field is truly
optional.  Required fields SHOULD use plain \texttt{\$\{order.total\}} so that
\texttt{GIS-PATH-MISSING} remains a safety net.  (Source: proposal \S5.5;
ratified decision rows ``Existing fixture migration'' and ``Hard-error default
unchanged''.)

% ─────────────────────────────────────────────────────────────────────────────
\subsection{Out of Scope}
\label{subsec:gis:optional-chaining:out-of-scope}

This section does not specify optional chaining for GXL boolean expressions;
\texttt{?.} is GIS-only in this extension.  It does not specify optional
chaining for GCP capture paths; use \texttt{capture.default:} at the capture
layer where that mechanism applies (\S\ref{sec:gcp:default}).  It does not add
nullish coalescing: \texttt{??} is deferred and MUST NOT be accepted as part of
this extension.  It also does not specify the audit trace event format for an
optional miss.  (Source: ratified decision rows ``Scope'', ``\texttt{??}
nullish-coalescing'', and ``Audit trace on optional miss''.)

\begin{quote}
\noindent\textit{Non-normative note.}
The proposed \texttt{GIS-PATH-OPTIONAL-MISS} info-level diagnostic is deferred
to the runtime design phase.  This chapter locks the syntax and evaluation
result; it does not define JSONL audit-trace event shape, payload fields, or
emission policy.
\end{quote}
